% GENERATED FILE. Edit canonical lessons, then run npm run generate:books.
% canonical-source-hash: fnv1a64:083074b5909079ec
% canonical-lessons: BN-C274-jodi-tahole, BN-C274-brishti-hole, BN-C274-dhoka, BN-W274-letter-ddha, BN-R274-if-and-otherwise

\chapter{If, and Otherwise}
\label{ch:bn-if-and-otherwise}
\hlchaptermodality{274}

\begin{chapteropening}
\textbf{By the end of this chapter:} \emph{I can say what will happen if (jôdi … tāhôle), say it with the -le form (brishṭi hôle), add otherwise (nāhôle), and read and write \bn{ঢ} in \bn{ঢোকা}.}

Complete BN-R274-if-and-otherwise: say what will happen if, and otherwise.
\end{chapteropening}

\section[\texorpdfstring{\bn{যদি} … \bn{তাহলে} … (jôdi … tāhôle …)}{যদি … তাহলে … (jôdi … tāhôle …)}]{\bn{যদি} … \bn{তাহলে} … (jôdi … tāhôle …) — if, then}

\label{lesson:BN-C274-jodi-tahole}



\begin{quote}
Say \textquotedblleft{}for two years\textquotedblright{}. (\emph{dui bôchhôr dhôre}.) Then \textbf{\bn{যদি}}, \textquotedblleft{}if\textquotedblright{}, and \textbf{\bn{তাহলে}}, \textquotedblleft{}then, in that case\textquotedblright{}: both of which you know.
\end{quote}

\subsection*{You'll want to know: if, then}
\begin{quote}
\textbf{\bn{যদি} \bn{কাল} \bn{ঝড়} \bn{হয়}, \bn{তাহলে} \bn{আমরা} \bn{যাব} \bn{না।}}
\end{quote}

\begin{quote}
\emph{jôdi kāl jhôṛ hôy, tāhôle āmrā jābô nā.} — \textquotedblleft{}If there's a storm tomorrow, we won't go.\textquotedblright{}
\end{quote}

The \textbf{\bn{যদি}} half takes the present, even about tomorrow: \textbf{\bn{হয়}}, not \textbf{\bn{হবে}}. The \textbf{\bn{তাহলে}} half takes the future. \textbf{\bn{তাহলে}} may be left out: \textbf{\bn{যদি} \bn{তুমি} \bn{আসো}, \bn{আমি} \bn{রান্না} \bn{করব}} (\emph{jôdi tumi āso, āmi rānnā kôrbô}), \textquotedblleft{}if you come, I'll cook\textquotedblright{}.

\subsection*{You'll want to know: if not}
Inside a \textbf{\bn{যদি}} half, \textbf{\bn{না}} comes before the verb, not after it:

\begin{quote}
\textbf{\bn{যদি} \bn{ট্রেন} \bn{সময়মতো} \bn{না} \bn{আসে}, \bn{তাহলে} \bn{আমি} \bn{বাসে} \bn{যাব।}}
\end{quote}

\begin{quote}
\emph{jôdi ṭren sômôymôto nā āse, tāhôle āmi bāse jābô.} — \textquotedblleft{}If the train doesn't come on time, I'll take the bus.\textquotedblright{}
\end{quote}

\subsection*{Your turn}
Say these aloud:

\begin{itemize}
\raggedright
  \item \textquotedblleft{}If there's a storm tomorrow, we won't go\textquotedblright{} — \emph{jôdi kāl jhôṛ hôy, tāhôle āmrā jābô nā}
  \item \textquotedblleft{}If you come, I'll cook\textquotedblright{} — \emph{jôdi tumi āso, āmi rānnā kôrbô}
  \item \textquotedblleft{}If the train doesn't come on time, I'll take the bus\textquotedblright{} — \emph{jôdi ṭren sômôymôto nā āse, tāhôle āmi bāse jābô}
\end{itemize}

\begin{tcolorbox}[breakable,colback=teal!4,colframe=teal!35!black,title={Before you move on}]
Which tense after \textbf{\bn{যদি}}? (The present.) After \textbf{\bn{তাহলে}}? (The future.) Where does \textbf{\bn{না}} go in the \textbf{\bn{যদি}} half? (Before the verb.)
\end{tcolorbox}
\section[\texorpdfstring{\bn{বৃষ্টি} \bn{হলে} … (brishṭi hôle …)}{বৃষ্টি হলে … (brishṭi hôle …)}]{\bn{বৃষ্টি} \bn{হলে} … (brishṭi hôle …) — the -le form, and otherwise}

\label{lesson:BN-C274-brishti-hole}



\begin{quote}
Say \textquotedblleft{}If there's a storm tomorrow, we won't go\textquotedblright{}. (\emph{jôdi kāl jhôṛ hôy, tāhôle āmrā jābô nā}.) Then \textbf{\bn{নাহলে}}, \textquotedblleft{}otherwise\textquotedblright{}.
\end{quote}

\subsection*{You'll want to know: the -le form}
Bengali can say \textquotedblleft{}if\textquotedblright{} with no \textbf{\bn{যদি}}: the verb takes \textbf{-\bn{লে}}, on the stem of the past, and stays the same for every person.

\begin{quote}
\textbf{\bn{বৃষ্টি} \bn{হলে} \bn{আমরা} \bn{বাড়িতে} \bn{থাকব।}}
\end{quote}

\begin{quote}
\emph{brishṭi hôle āmrā bāṛite thākbô.} — \textquotedblleft{}If it rains, we'll stay at home.\textquotedblright{}
\end{quote}

\textbf{\bn{বৃষ্টি}} (\emph{brishṭi}) is \textquotedblleft{}rain\textquotedblright{}, and \textbf{\bn{বৃষ্টি} \bn{হওয়া}} is \textquotedblleft{}to rain\textquotedblright{}: \textquotedblleft{}rain happens\textquotedblright{}.

\textbf{\bn{হওয়া}} $\to$ \textbf{\bn{হলে}}, \textbf{\bn{আসা}} $\to$ \textbf{\bn{এলে}}, \textbf{\bn{যাওয়া}} $\to$ \textbf{\bn{গেলে}}, \textbf{\bn{খাওয়া}} $\to$ \textbf{\bn{খেলে}}, \textbf{\bn{করা}} $\to$ \textbf{\bn{করলে}}, \textbf{\bn{কেনা}} $\to$ \textbf{\bn{কিনলে}}: \textbf{\bn{তুমি} \bn{এলে} \bn{আমি} \bn{খুশি} \bn{হব}} (\emph{tumi ele āmi khushi hôbô}), \textquotedblleft{}if you come, I'll be happy\textquotedblright{}. The person is in the noun before it, never in the ending.

\subsection*{You'll want to know: otherwise}
\textbf{\bn{নাহলে}}, \textquotedblleft{}otherwise\textquotedblright{}, which you know, is \textbf{\bn{না}} with this same \textbf{\bn{হলে}}: \textquotedblleft{}if it is not so\textquotedblright{}.

\begin{quote}
\textbf{\bn{তাড়াতাড়ি} \bn{যেতে} \bn{হবে}, \bn{নাহলে} \bn{ট্রেন} \bn{ছেড়ে} \bn{যাবে।}}
\end{quote}

\begin{quote}
\emph{tāṛātāṛi jete hôbe, nāhôle ṭren chheṛe jābe.} — \textquotedblleft{}We have to go quickly, otherwise the train will leave.\textquotedblright{}
\end{quote}

\subsection*{Your turn}
Say these aloud:

\begin{itemize}
\raggedright
  \item \textquotedblleft{}If it rains, we'll stay at home\textquotedblright{} — \emph{brishṭi hôle āmrā bāṛite thākbô}
  \item \textquotedblleft{}If you come, I'll be happy\textquotedblright{} — \emph{tumi ele āmi khushi hôbô}
  \item \textquotedblleft{}We have to go quickly, otherwise the train will leave\textquotedblright{} — \emph{tāṛātāṛi jete hôbe, nāhôle ṭren chheṛe jābe}
\end{itemize}

\begin{tcolorbox}[breakable,colback=teal!4,colframe=teal!35!black,title={Before you move on}]
\textbf{\bn{আসা}} in the \textbf{-\bn{লে}} form? (\textbf{\bn{এলে}}.) Does \textbf{-\bn{লে}} change for the person? (No.) \textquotedblleft{}Otherwise\textquotedblright{}? (\textbf{\bn{নাহলে}}: \textquotedblleft{}if not so\textquotedblright{}.)
\end{tcolorbox}
\section[\texorpdfstring{\bn{ঢোকা} (ḍhokā)}{ঢোকা (ḍhokā)}]{\bn{ঢোকা} (ḍhokā) — to enter}

\label{lesson:BN-C274-dhoka}



\begin{quote}
Before the new one: say \textquotedblleft{}If you come, I'll be happy\textquotedblright{}. (\emph{tumi ele āmi khushi hôbô}.)
\end{quote}

\subsection*{You'll want to know: \bn{ঢোকা}}
\textbf{\bn{ঢোকা}} — \emph{ḍhokā} — \textquotedblleft{}to enter, to go in\textquotedblright{}.

\begin{quote}
\textbf{\bn{বাড়িতে} \bn{ঢুকে} \bn{আমি} \bn{জুতো} \bn{খুললাম।}}
\end{quote}

\begin{quote}
\emph{bāṛite ḍhuke āmi juto khullām.} — \textquotedblleft{}I came into the house and took my shoes off.\textquotedblright{}
\end{quote}

The place you enter takes \textbf{-\bn{এ}} or \textbf{-\bn{তে}}: \textbf{\bn{বাড়িতে}}, \textbf{\bn{দোকানে}}. The \textbf{\bn{ও}} of \textbf{\bn{ঢোকা}} rises to \textbf{\bn{উ}} in the past, the \emph{-e} form and the \textbf{-\bn{লে}} form, as \textbf{\bn{খোলা}} does: \textbf{\bn{সে} \bn{দোকানে} \bn{ঢুকল}} (\emph{she dokāne ḍhuklô}), \textquotedblleft{}she went into the shop\textquotedblright{}; \textbf{\bn{ঢুকে}}; \textbf{\bn{ঢুকলে}}.

\begin{cousinweb}
\textbf{\bn{ঢোকা}} (\emph{ḍhokā}) is listed in CDIAL, as \emph{ḍhukā}, \textquotedblleft{}to enter\textquotedblright{}, under a reconstructed \emph{*ḍhukyati}, \textquotedblleft{}approaches\textquotedblright{} (CDIAL 5592): a form no text records, hence the star.
\end{cousinweb}

\subsection*{Your turn}
Say these aloud:

\begin{itemize}
\raggedright
  \item \emph{ḍhokā}, then \emph{she ḍhuklô}
  \item \textquotedblleft{}I came into the house and took my shoes off\textquotedblright{} — \emph{bāṛite ḍhuke āmi juto khullām}
  \item \textquotedblleft{}She went into the shop\textquotedblright{} — \emph{she dokāne ḍhuklô}
\end{itemize}

\begin{tcolorbox}[breakable,colback=teal!4,colframe=teal!35!black,title={Before you move on}]
What does \textbf{\bn{ঢোকা}} mean? (\textquotedblleft{}To enter\textquotedblright{}.) Its past? (\textbf{\bn{ঢুকল}}, with \textbf{\bn{উ}}.) Its \emph{-e} form? (\textbf{\bn{ঢুকে}}.)
\end{tcolorbox}
\section[\texorpdfstring{\bn{ঢ} (ḍhô)}{ঢ (ḍhô)}]{\bn{ঢ} (ḍhô) — the letter at the start of \bn{ঢোকা}}

\label{lesson:BN-W274-letter-ddha}



\begin{quote}
Before the new one: say \textbf{\bn{ঢোকা}}, \textquotedblleft{}to enter\textquotedblright{}. Then read \textbf{\bn{ট}} (\emph{ṭô}) and \textbf{\bn{ড}} (\emph{ḍô}).
\end{quote}

\subsection*{Script you'll notice: \bn{ঢ}}
\textbf{\bn{ঢ}} — \emph{ḍhô}: the \textbf{\bn{ড}} sound with a breath after it, as \textbf{\bn{ঠ}} is \textbf{\bn{ট}} with a breath. It opens \textbf{\bn{ঢোকা}}.

In the print of this book, \textbf{\bn{ঢ}} is \textbf{\bn{ট}} without its raised hook. It has a headline; a straight stem stands under the left part of it; the stem's foot rounds into a bowl, and the bowl's right side ends in a solid knob, tucked inside the curve under the right half of the headline. The right side does not reach the headline. Look at \textbf{\bn{ট}} beside it: below the headline the two are alike, and \textbf{\bn{ট}} has the hook above.

\subsection*{Writing: \bn{ঢ} — copy the printed shape}
Put your pen beside \textbf{\bn{ঢ}} and copy the shape you can see, slowly and large. Check the copy against the print: the headline, the stem on the left, the round bottom, the solid knob, and the gap between the knob and the headline.

\begin{quote}
This book does not yet tell you where to start \bn{ঢ} or which way the pen goes. That is not written down here until it can be written down with a source. Copying the printed shape needs no such source, and it is how the shape becomes yours in the meantime.
\end{quote}

\subsection*{Your turn}
\begin{itemize}
\raggedright
  \item \emph{Find:} \bn{ঢ} at the front of \textbf{\bn{ঢোকা}}
  \item \emph{Write it:} \bn{ট}, then \bn{ঢ}, from the print
  \item \emph{Write it:} \textbf{\bn{ঢোকা}}, then \textbf{\bn{ঢুকল}}
  \item \emph{Write it:} \textbf{\bn{দৈনিক}}, from memory
\end{itemize}

\begin{tcolorbox}[breakable,colback=teal!4,colframe=teal!35!black,title={Before you move on}]
Which letter opens \textbf{\bn{ঢোকা}}? (\textbf{\bn{ঢ}}.) What does its right side end in? (A solid knob.) What does \textbf{\bn{ট}} have that \textbf{\bn{ঢ}} does not? (The hook above the headline.)
\end{tcolorbox}
\section[(review)]{Review — if, and otherwise}

\label{lesson:BN-R274-if-and-otherwise}



\begin{quote}
Say \textquotedblleft{}If there's a storm tomorrow, we won't go\textquotedblright{}, then \textquotedblleft{}if it rains tomorrow, we won't go\textquotedblright{}, with \textbf{\bn{হলে}}. (\emph{jôdi kāl jhôṛ hôy, tāhôle āmrā jābô nā}; \emph{kāl brishṭi hôle āmrā jābô nā}.)
\end{quote}

\subsection*{Your turn}
\begin{itemize}
\raggedright
  \item \emph{Say it:} If you have a fever, you ought to go to the doctor.
  \item \emph{Say it:} If the bus doesn't come, we'll walk.
  \item \emph{Say it:} If Mira comes, we'll go to the market.
  \item \emph{Say it:} We have to leave now, otherwise we'll be late.
  \item \emph{Say it:} The shop is open from morning until night. She went in and bought a book.
  \item \emph{Write it:} \textbf{\bn{ঢোকা}}, then \textbf{\bn{দৈনিক}}
\end{itemize}

\begin{tcolorbox}[breakable,colback=teal!4,colframe=teal!35!black,title={Before you move on}]
Which tense after \textbf{\bn{যদি}}? (The present.) \textbf{\bn{যাওয়া}} in the \textbf{-\bn{লে}} form? (\textbf{\bn{গেলে}}.) \textquotedblleft{}Otherwise\textquotedblright{}? (\textbf{\bn{নাহলে}}.) Which letter opens \textbf{\bn{ঢোকা}}? (\textbf{\bn{ঢ}}.)
\end{tcolorbox}
